% Mathematical proof material for the main gamma/iPRA article.
% Dependencies: elementary exp and positive-real log with finite error bounds.
% No global complex logarithm, analytic continuation theorem, or choice is used.

\section{The reciprocal function from finite Euler products}
\label{sec:reciprocal-products}

All bounds in this section are uniform in a displayed integer disk radius
$R\geq1$.  Bounds on complex absolute values can be implemented using rational
squared absolute values; an upper bound in the rational norm
$|\Re z|+|\Im z|$ is also sufficient.  Expressions involving nonrational
arguments mean the corresponding fixed primitive recursive names with their
given error certificates.  They do not introduce a quantifier over a universe
of all primitive recursive functions.

\subsection{Euler's constant with a modulus}
Put $H_N=\sum_{k=1}^N1/k$ and $\gamma_N=H_N-\log N$ for $N\geq1$.
The elementary integral or logarithmic-series inequalities give
\[
 0\leq\log(1+1/k)-\frac1{k+1}
 \leq\frac1k-\frac1{k+1}.
\]
Consequently, for $M\geq N$,
\[
 0\leq\gamma_N-\gamma_M\leq\frac1N-\frac1M\leq\frac1N.
 \tag{C1}
\]
Moreover $0\leq\gamma_N\leq1$, by comparing the finite harmonic sum with
the integral of $1/t$ over $[1,N]$.  Thus (C1), together with a finite
approximation of $\log N$, defines a regular name for $\gamma$, and
\[
 0\leq\gamma_N-\gamma\leq N^{-1},\qquad 0\leq\gamma\leq1.
 \tag{C2}
\]
For example $N=2^{p+2}$ and a logarithm evaluation within $2^{-p-2}$
suffice for an approximation with error less than $2^{-p-1}$.
Monotone convergence has not been invoked: (C1) is the Cauchy certificate.

\subsection{A local logarithm that has no branch ambiguity}
For $|w|\leq1/2$ let
\[
 L(w)=\sum_{j=2}^{\infty}\frac{(-1)^{j+1}w^j}{j}.
\]
The series has the explicit estimates
\[
 |L(w)|\leq\frac{|w|^2}{2(1-|w|)}\leq|w|^2,
 \qquad
 \left|L(w)-\sum_{j=2}^{d}\frac{(-1)^{j+1}w^j}{j}\right|
 \leq\frac{|w|^{d+1}}{(d+1)(1-|w|)}.
 \tag{C3}
\]
Here $w+L(w)$ is only the power-series logarithm near $1$.
The identity
\[
 \exp\bigl(w+L(w)\bigr)=1+w
 \tag{C4}
\]
follows directly by finite coefficient algebra.  Write
$F(w)=\exp(w+L(w))=\sum_{n\geq0}f_nw^n$.
The geometric derivative of $w+L(w)$ gives $(1+w)F'=F$ and $f_0=1$.
Comparing finite coefficients yields $f_1=1$ and
$(n+1)f_{n+1}=(1-n)f_n$ for $n\geq1$; hence every $f_n$ with
$n\geq2$ is zero.  Differentiation and multiplication are verified first
for finite polynomials; their geometric and exponential tails bound each
discarded term.  This uses neither a global logarithm on $\mathbb C$
nor an assertion about branches.

\subsection{The entire reciprocal and quantitative Euler convergence}
Define finite entire expressions
\[
 Q_N(z)=z\exp(-z\log N)\prod_{k=1}^N(1+z/k),\qquad
 P_N(z)=z\exp(\gamma z)\prod_{k=1}^N(1+z/k)\exp(-z/k).
 \tag{C5}
\]
Fix $|z|\leq R$ and set $K=2R$.  For $N\geq K$ the tail series
\[
 T_N(z)=\sum_{k=N+1}^{\infty}L(z/k)
\]
has the finite-tail certificate
\[
 |T_N(z)|\leq R^2\sum_{k>N}k^{-2}\leq\frac{R^2}{N}.
 \tag{C6}
\]
The last inequality follows already for every finite sum by
$k^{-2}\leq1/[k(k-1)]$, so that it is a telescoping rational inequality.
Define
\[
 P(z)=z\prod_{k=1}^K(1+z/k)
       \exp\bigl((\gamma-H_K)z+T_K(z)\bigr).
 \tag{C7}
\]
The definitions with different sufficiently large $K$ agree, by (C4) and
finite product cancellation.  In particular
\[
 P(z)=P_N(z)\exp(T_N(z)),\qquad
 Q_N(z)=P_N(z)\exp\bigl(z(\gamma_N-\gamma)\bigr).
 \tag{C8}
\]
The quotients suggested by this notation are never needed at a zero: (C8)
is an equality obtained by multiplication, valid everywhere.

For an entirely rational majorant put
\[
 M_R=R(1+R)^{2R}3^{2R^2+2R},\qquad B_R=R^2+R.
 \tag{C9}
\]
Indeed, the absolute value of the polynomial in (C7) is at most
$R(1+R)^{2R}$, while the exponent has real part at most
\[
 (1+H_{2R})R+R^2/(2R)
 \leq (1+2R)R+R/2\leq2R^2+2R.
\]
Using $e\leq3$ proves $|P(z)|\leq M_R$, and the same estimate bounds
$P_N(z)$ for $N\geq2R$.  By (C2), (C6), (C8), and
$|e^u-e^v|\leq e^{|v|}e^{|u-v|}|u-v|$, or directly by the
exponential power series, one obtains
\[
 |Q_N(z)-P(z)|\leq \frac{3M_RB_R}{N}
 \quad\text{if }N\geq\max(2R,B_R).
 \tag{C10}
\]
One convenient derivation writes both sides using (C7); the difference
of their tail exponents has absolute value at most $B_R/N\leq1$.
Therefore $N\geq3M_RB_R2^p$ gives an error at most $2^{-p}$.
Evaluating the finitely many exponentials and logarithms with a further
error allowance gives a rational algorithm, since all the displayed
lengths and precision budgets are primitive recursive.

Thus (C10) constructs $P$ on every disk.  Equations (C5)--(C8) also prove
the reciprocal Weierstrass product
\[
 \boxed{\displaystyle
 P(z)=z e^{\gamma z}\prod_{k=1}^{\infty}(1+z/k)e^{-z/k}.}
 \tag{C11}
\]
The infinite product means precisely its explicitly controlled finite
products.  It has not been inferred from an abstract product convergence
theorem.

\subsection{An explicit lower bound away from the poles}
Suppose, in addition to $|z|\leq R$, that a rational $0<\delta\leq1$
satisfies
\[
 |z+k|\geq\delta\quad(0\leq k\leq2R).
 \tag{C12}
\]
Set
\[
 \lambda_{R,\delta}=
 \frac{\delta^{2R+1}}{(2R)!\,3^{2R^2+2R}}.
 \tag{C13}
\]
In (C7), the polynomial has absolute value at least
$\delta^{2R+1}/(2R)!$.  The real part of its exponent is at least
$-(2R^2+2R)$, so
\[
 |P(z)|\geq\lambda_{R,\delta}>0.
 \tag{C14}
\]
The identical argument with $\gamma_N$ and the finite tail instead of
$\gamma$ and $T_K$ gives $|Q_N(z)|\geq\lambda_{R,\delta}$ for
$N\geq2R$.  This is a positive lower bound, not merely the negation of an
equality.  It therefore licenses constructive inversion.

On a disk supplied with (C12) define $\Gamma(z)=1/P(z)$.  Its definition
agrees on overlaps.  Its Euler approximants are
\[
 \Gamma_N(z)=\frac{N!N^z}{z(z+1)\cdots(z+N)}=Q_N(z)^{-1}.
\]
For $N\geq\max(2R,B_R)$, (C10)--(C14) yield
\[
 \boxed{\displaystyle
 |\Gamma_N(z)-\Gamma(z)|
 \leq\frac{3M_RB_R}{\lambda_{R,\delta}^{\,2}N}.}
 \tag{C15}
\]
Also $|\Gamma(z)|\leq1/\lambda_{R,\delta}$ and
$|\Gamma(z)|\geq1/M_R$.  In particular gamma has no zeros, with a
quantitative meaning on each supplied disk.  When $\Re z\geq a>0$,
one may take $\delta=\min(a,1)$; hence (C15) identifies this construction
with any independently proved Euler limit on the right half-plane,
including the Euler integral construction.

\subsection{Why ``entire'' is justified constructively}
We give the missing local analytic information rather than applying a
classical theorem about uniformly convergent holomorphic functions.
For $K=2R$, finite rearrangement with the absolute bounds (C3) gives
\[
 T_K(z)=\sum_{m=2}^{\infty}
 \frac{(-1)^{m+1}z^m}{m}
 \left(\zeta(m)-H_K^{(m)}\right),\qquad
 H_K^{(m)}=\sum_{k=1}^K k^{-m}.
 \tag{C16}
\]
Here the positive integer zeta values are defined by their own rational
partial sums, with
\[
 0\leq\zeta(m)-\sum_{k=1}^N k^{-m}
 \leq\frac{1}{(m-1)N^{m-1}}\quad(m\geq2).
 \tag{C17}
\]
The integral comparison in (C17) can equally be replaced by its finite
antiderivative inequality and rational telescoping estimates.
In (C16) the tail after degree $d\geq1$ is at most
\[
 \frac{R^2}{K}\,2^{1-d}.
 \tag{C18}
\]
For example, its $m$th term is bounded by
$(R^2/K)2^{-(m-2)}/[m(m-1)]$, whose geometric tail proves (C18).

Finite polynomial multiplication and the exponential power series now
turn (C7) into a convergent power series on each disk.  This operation has
a direct absolute majorant: if $\sum_j|a_j|t^j\leq A$, then the
coefficient absolute sum of $\exp(\sum_j a_jz^j)$ at $t$ is at most
$e^A$.  This is proved by finite Cauchy multiplication and the ordinary
exponential tail estimate.  There is no appeal to an infinite formal
manipulation without an error bound.

For a useful global consequence, make the construction at radius $2R$
and use $K=4R$.  Its coefficient absolute sum at $2R$ is bounded by
$M_{2R}$.  Consequently, writing $P(z)=\sum_{j\geq1}c_jz^j$,
\[
 \left|P(z)-\sum_{j=1}^d c_jz^j\right|
 \leq M_{2R}2^{-d-1}\quad(|z|\leq R).
 \tag{C19}
\]
The coefficients do not depend on the auxiliary $K$: equality of each
finite-degree coefficient follows by expanding (C4) and cancelling one
factor at a time.  Coefficients are thus computable by finite recursive
convolutions with the constants in (C2) and (C17).

For completeness, a power series about $z_0$, whose coefficient absolute
sum at radius $r$ is at most $A$, satisfies
\[
 |f(z_0+h)-f(z_0)-b_1h|\leq A|h|^2/r^2
 \quad(|h|\leq r),
\]
when $b_1$ is its first coefficient about $z_0$.
The series about any $z_0$ is obtained by the binomial formula, using a
larger disk to bound all rearrangement tails.  The displayed remainder
is the definition and a modulus of its complex derivative.  Repeated
differentiation is handled with an extra disk margin.  This proves
entireness in the concrete sense needed here.  For inversion, suppose
$|P(z_0)|\geq d>0$ and the coefficient absolute sum about $z_0$ is at
most $A$ at radius $r$.  Then $A\geq d>0$.  Choose a positive rational
\[
 \rho\leq\min\left(\frac r2,\frac{dr}{4A}\right).
\]
On $|h|\leq\rho$, the nonconstant coefficient tail gives
$|P(z_0+h)-P(z_0)|\leq A\rho/r\leq d/4$.
The geometric inverse therefore has ratio at most $1/4$ (and, in
particular, at most $1/2$), which provides its convergence and remainder
modulus.  This proves analyticity of $1/P$ on the smaller disk.  Thus
meromorphic language below is supported by local power series and
explicit division bounds.

\section{Functional equation, zeros, poles, and conjugation}

Finite cancellation, without division by $z$, gives
\[
 zQ_N(z+1)=\left(1+\frac{z+1}{N}\right)Q_N(z).
\]
The bounds (C10) give
\[
 P(z)=zP(z+1),\qquad P(1)=1,
 \tag{C20}
\]
the latter also following from $Q_N(1)=(N+1)/N$.  Therefore, wherever
the divisions are licensed,
\[
 \Gamma(z+1)=z\Gamma(z),\qquad\Gamma(1)=1,\qquad
 \Gamma(n+1)=n!.
 \tag{C21}
\]
The last assertion uses ordinary finite induction on $n$; its analytic
content has already been reduced to (C20).

Iteration of (C20) gives the entire identity
\[
 P(z)=\prod_{j=0}^{n}(z+j)\,P(z+n+1).
 \tag{C22}
\]
At $z=-n$ this proves $P(-n)=0$, and differentiating the finite product
proves
\[
 P'(-n)=(-1)^n n!.
 \tag{C23}
\]
So all these zeros are simple.

The converse zero assertion can also be made constructively; it does
not require the invalid general rule that a zero product of arbitrary
constructive reals has a zero factor.  Suppose $P(z)=0$ and $|z|\leq R$.
Compute a rational complex $q$ with $|q-z|\leq1/16$.  Decide the finite
list of rational tests $|q+j|^2<1/16$ for $0\leq j\leq2R$.
If all fail, then $|z+j|\geq3/16$ for all $j$, contradicting (C14).
Hence one test succeeds.  There is at most one such $j$, because two
different integer centers have distance at least $1$.  For this $j$,
$|z+j|<5/16$, while $|z+k|>11/16$ for every other integer $k$.
Factor $(z+j)$ out of (C7).  The remaining factor has a positive lower
bound by exactly the proof of (C14), using $11/16$ for its finitely many
linear factors.  Hence $P(z)=0$ implies $z=-j$.  The bounded search for
$j$ uses a decidable rational predicate and so is primitive recursive.

Equation (C22) gives a local description of gamma at every pole:
\[
 \Gamma(-n+h)=
 \frac{(-1)^n}{n!\,h}\,
 \Gamma(1+h)\prod_{j=1}^n(1-h/j)^{-1}.
 \tag{C24}
\]
All factors apart from $h^{-1}$ have convergent power series for
$|h|<1$.  Their value at zero is $1$, proving a simple pole with residue
\[
 \operatorname{Res}_{z=-n}\Gamma(z)=\frac{(-1)^n}{n!}.
 \tag{C25}
\]
No value is assigned to gamma at its poles.  A totalized implementation
may choose a default value there, but the meromorphic assertions concern
the punctured local formula (C24).

All coefficients and constants in (C5) are real.  Conjugating finite
expressions and using (C10), (C14) proves
\[
 P(\overline z)=\overline{P(z)},\qquad
 \Gamma(\overline z)=\overline{\Gamma(z)}.
 \tag{C26}
\]

